\chapter{Аналитическая часть}
В данной части работы будут рассмотрены бинарный алгоритм поиска и алгоритм поиска полным перебором.

\section{Алгоритм поиска полным перебором}
Идея данного алгоритма заключается в последовательном рассмотрении каждого элемента массива и сравнении его с искомым значением до тех пор, пока это сравнение не окажется истинным, или в массиве не закончатся элементы. Таким образом в случае, когда в массиве нет искомого элемента, необходимо провести N сравнений. Однако, если искомый элемент находится в массиве на последнем месте, чтобы его найти нужно также провести N сравнений. Значит, худший случай наступает в двух возможных ситуациях, а алгоритм имеет асимптотическую оценку $O(n)$.

\section{Алгоритм бинарного поиска}
Данный алгоритм предполагает, что на вход данные подаются в упорядоченном виде.
Изначально необходимо сравнить искомый элемент $x$ с элементом массива, находящимся в середине. На основе результата этого сравнения определяется та половина массива, в которой необходимо продолжить поиск. Далее действия выполняются аналогично, но уже к определённому на предыдущем шаге подмассиву, до тех пор, пока индекс левой части рассматриваемого подмассива не станет больше индекса правой части подмассива, или пока искомый элемент не будет найден. Таким образом, для установления позиции искомого элемента или факта его отсутствия в массиве необходимо провести не более, чем $log_{2}{N}$ сравнений. Такая процедура иногда называется <<логарифмическим поиском>> или <<методом деления пополам>>, но наиболее употребительный термин~---~бинарный поиск~\cite{book_knut}. Для полной оценки трудоёмкости алгоритма необходимо также учесть трудоёмкость используемой сортировки, поскольку данные в задаче могут быть поданы не в упорядоченном виде. 

\section*{Вывод}
В данной части были кратко рассмотрены бинарный алгоритм поиска и алгоритм поиска полным перебором.
